% Supplementary A — Journal paper §I (Introduction)
% Extended motivation, historical context, problem framing.

\documentclass[11pt,a4paper]{article}
% Shared preamble for all supplementary chapter documents.
% Used by each supplementary/conference/suppl_*.tex and
% supplementary/journal/suppl_*.tex via % Shared preamble for all supplementary chapter documents.
% Used by each supplementary/conference/suppl_*.tex and
% supplementary/journal/suppl_*.tex via % Shared preamble for all supplementary chapter documents.
% Used by each supplementary/conference/suppl_*.tex and
% supplementary/journal/suppl_*.tex via \input{../preamble.tex}.

\usepackage[utf8]{inputenc}
\usepackage[T1]{fontenc}
\usepackage[a4paper,margin=2.2cm]{geometry}
\usepackage{amsmath,amssymb,amsfonts,amsthm}
\usepackage{mathtools}
\usepackage{graphicx}
\usepackage{booktabs}
\usepackage{array}
\usepackage{multirow}
\usepackage{longtable}
\usepackage{algorithm}
\usepackage{algpseudocode}
\usepackage{xcolor}
\usepackage{url}
\usepackage{cite}
\usepackage[hidelinks]{hyperref}

\graphicspath{{../../figures/}}

\renewcommand{\thesection}{\Alph{section}.\arabic{section}}
\renewcommand{\thesubsection}{\thesection.\arabic{subsection}}

\theoremstyle{plain}
\newtheorem{theorem}{Theorem}[section]
\newtheorem{lemma}[theorem]{Lemma}
\newtheorem{proposition}[theorem]{Proposition}
\theoremstyle{definition}
\newtheorem{definition}[theorem]{Definition}
\theoremstyle{remark}
\newtheorem{remark}[theorem]{Remark}

\setcounter{secnumdepth}{3}
\setcounter{tocdepth}{2}

\newcommand{\R}{\mathbb{R}}
\newcommand{\E}{\mathbb{E}}
\newcommand{\KL}[2]{\mathrm{KL}\!\left(#1 \,\middle\|\, #2\right)}
\newcommand{\ari}{\mathrm{ARI}}
\newcommand{\nmi}{\mathrm{NMI}}
\newcommand{\softmax}{\mathrm{softmax}}
.

\usepackage[utf8]{inputenc}
\usepackage[T1]{fontenc}
\usepackage[a4paper,margin=2.2cm]{geometry}
\usepackage{amsmath,amssymb,amsfonts,amsthm}
\usepackage{mathtools}
\usepackage{graphicx}
\usepackage{booktabs}
\usepackage{array}
\usepackage{multirow}
\usepackage{longtable}
\usepackage{algorithm}
\usepackage{algpseudocode}
\usepackage{xcolor}
\usepackage{url}
\usepackage{cite}
\usepackage[hidelinks]{hyperref}

\graphicspath{{../../figures/}}

\renewcommand{\thesection}{\Alph{section}.\arabic{section}}
\renewcommand{\thesubsection}{\thesection.\arabic{subsection}}

\theoremstyle{plain}
\newtheorem{theorem}{Theorem}[section]
\newtheorem{lemma}[theorem]{Lemma}
\newtheorem{proposition}[theorem]{Proposition}
\theoremstyle{definition}
\newtheorem{definition}[theorem]{Definition}
\theoremstyle{remark}
\newtheorem{remark}[theorem]{Remark}

\setcounter{secnumdepth}{3}
\setcounter{tocdepth}{2}

\newcommand{\R}{\mathbb{R}}
\newcommand{\E}{\mathbb{E}}
\newcommand{\KL}[2]{\mathrm{KL}\!\left(#1 \,\middle\|\, #2\right)}
\newcommand{\ari}{\mathrm{ARI}}
\newcommand{\nmi}{\mathrm{NMI}}
\newcommand{\softmax}{\mathrm{softmax}}
.

\usepackage[utf8]{inputenc}
\usepackage[T1]{fontenc}
\usepackage[a4paper,margin=2.2cm]{geometry}
\usepackage{amsmath,amssymb,amsfonts,amsthm}
\usepackage{mathtools}
\usepackage{graphicx}
\usepackage{booktabs}
\usepackage{array}
\usepackage{multirow}
\usepackage{longtable}
\usepackage{algorithm}
\usepackage{algpseudocode}
\usepackage{xcolor}
\usepackage{url}
\usepackage{cite}
\usepackage[hidelinks]{hyperref}

\graphicspath{{../../figures/}}

\renewcommand{\thesection}{\Alph{section}.\arabic{section}}
\renewcommand{\thesubsection}{\thesection.\arabic{subsection}}

\theoremstyle{plain}
\newtheorem{theorem}{Theorem}[section]
\newtheorem{lemma}[theorem]{Lemma}
\newtheorem{proposition}[theorem]{Proposition}
\theoremstyle{definition}
\newtheorem{definition}[theorem]{Definition}
\theoremstyle{remark}
\newtheorem{remark}[theorem]{Remark}

\setcounter{secnumdepth}{3}
\setcounter{tocdepth}{2}

\newcommand{\R}{\mathbb{R}}
\newcommand{\E}{\mathbb{E}}
\newcommand{\KL}[2]{\mathrm{KL}\!\left(#1 \,\middle\|\, #2\right)}
\newcommand{\ari}{\mathrm{ARI}}
\newcommand{\nmi}{\mathrm{NMI}}
\newcommand{\softmax}{\mathrm{softmax}}


\title{Journal Supplementary A --- The accuracy-only view and why it fails\\
{\large Companion to: \emph{Beyond Accuracy: A Multi-Axis Evaluation
Framework for Interpretable Topic Models on HSI}}}
\author{Felipe Santib\'a\~nez-Leal}
\date{2026}

\begin{document}
\maketitle

\section{Scope}
This supplement expands Section~I of the journal paper. We provide
(a)~a longer historical account of how the HSI community arrived at
the accuracy-only evaluation discipline, (b)~five concrete examples
in which a single OA number hides materially different model
behaviours, (c)~the philosophical commitment that distinguishes
\emph{interpretability} from \emph{post-hoc explanation}, and
(d)~the framing that motivates a twelve-axis framework rather than
a smaller or larger one.

\section{How the field arrived at OA-as-summary}

The earliest HSI classification papers
\cite{plaza2009recent,camps2014advances} settled on overall accuracy
(OA) and kappa as the canonical summary because the underlying task
was assumed to be \emph{closed-set supervised classification}: given
a labelled benchmark with $L$ classes, predict the class of each
unlabelled pixel; report the fraction correct. Under this
assumption OA is monotone in the quantity the practitioner cares
about (predictive accuracy) and the kappa adjustment handles class
imbalance. Subsequent work imported the OA discipline wholesale into
\emph{unsupervised} and \emph{interpretability-motivated} settings
where the quantity the practitioner cares about is different.

Linear unmixing~\cite{nascimento2005vca,bioucasdias2012unmixing}
imported OA-style reporting via downstream classification of
predicted abundance vectors. Deep classification papers
\cite{chen2014deep,paoletti2019deep} continued the OA tradition.
Topic-model adaptations to HSI~\cite{zou2017pmlda,liu2019hsilda}
inherited OA from both ancestors. The price of this inheritance is
that the OA number cannot itself answer the question
\emph{``why is this topic basis interpretable?''} --- because the
question is not what OA measures.

\section{Five concrete examples of OA hiding bad behaviour}

\paragraph{Example 1: seed-collapsed ProdLDA.} ProdLDA fits on
Indian Pines under five seeds yield ARI standard deviation $0.003$
across seeds (extremely stable) and per-seed ARI \emph{below} 0.17
(weak class discrimination). A single ProdLDA OA number averaged
over seeds would be honest but uninformative; the practitioner needs
to see both the stability and the level.

\paragraph{Example 2: high-OA on weakly identified topics.}
KSC under a 13-class supervised baseline reaches OA $> 0.9$ on the
canonical fold; the same KSC under our band-mask diagnostic returns
paired ARI $\approx 0.01$. A KSC paper that reports only OA would
\emph{look} like a strong contribution; the topic basis driving the
classification is, however, band-fragile.

\paragraph{Example 3: HIDSAG inversion.} On the labelled scenes the
topic-routed-soft classifier is statistically indistinguishable from
the raw-spectrum logistic baseline. On HIDSAG the topic-logistic
classifier sits at posterior $\mu = 0.316$ vs raw $\mu = 0.585$ with
$P[\mu_{\text{topic}} > \mu_{\text{raw}}] = 0.0$. The labelled-scene
OA number would lead a practitioner to deploy the topic
representation on HIDSAG; the hierarchical model says that decision
is wrong.

\paragraph{Example 4: coherence-vs-ARI flip.} ProdLDA's
top-$15$-word coherence is the highest on Indian Pines
($c_v = 0.630$), but its cluster-vs-label ARI is the lowest
($0.165$). A practitioner who picks ProdLDA on the basis of
coherence gets a different downstream behaviour than one who picks
on the basis of clustering quality.

\paragraph{Example 5: SLIC vs topic-dominant cross-method
disagreement.} On Indian Pines SLIC-500 and the canonical topic
dominant map have pairwise ARI of $0.039$ --- i.e.\ the two
``unsupervised partitions'' of the same scene disagree almost
completely on which pixels belong to which segment. The OA of
either method against ground truth does not surface this
disagreement; the cross-method ARI matrix does.

\section{Interpretability versus post-hoc explanation}

There is a methodological distinction we want to preserve.
\emph{Interpretability} is a property of the model itself: a topic
basis $\phi$ is interpretable if every element of the
parameterisation has a meaning the practitioner can read off the
parameter. \emph{Post-hoc explanation} is a property of an attached
diagnostic: a feature-attribution map or an inverse-classification
view that approximates the model's behaviour on a sample. Topic
models on HSI are advertised as belonging to the first category, but
in practice the published evaluations support only the second:
a high-OA topic model with a band-fragile basis is providing
post-hoc explanation, not interpretability.

The twelve-axis framework defended in the journal paper restores the
first-category claim or denies it explicitly. A topic basis whose
twelve-axis card has stable, coherence-meaningful, band-robust, and
sample-stable rows admits the
\emph{``$\phi_k$ has a material-level meaning''} reading. A topic
basis whose card has any of those rows degenerate admits only the
\emph{``$\theta$ is a useful feature for a downstream classifier''}
reading --- which is the post-hoc reading and is allowed, but should
be advertised as such.

\section{Why twelve axes, not five or fifty}

The number twelve is not magic. It reflects the orthogonalisation
of all the evaluation questions the present project has needed to
ask over $\sim 150$ iterative cycles of working with the artefact
set. We checked four alternative counts:

\begin{itemize}
\item \textbf{Five axes (the OA + four corollaries view).}
Insufficient. Drops cross-method agreement, topic--label coupling,
preprocessing robustness, cross-scene transfer.
\item \textbf{Seven axes.} Drops preprocessing robustness (essential
on HIDSAG) and cross-scene transfer (essential when sensor families
overlap, e.g.\ AVIRIS Salinas / Salinas-A).
\item \textbf{Twelve axes (selected).} The smallest set that yields a
distinct numeric per axis on the present benchmarks and admits a
deterministic builder per axis without overlapping builders.
\item \textbf{Twenty-plus axes.} We considered splitting B-3 into
inferential-stability + estimation-stability, B-5 into mask-set vs
band-budget, B-12 into intra-paper vs inter-paper baselines. Doing so
multiplies the artefact footprint without adding qualitatively new
information.
\end{itemize}

The twelve-axis count is therefore the minimum that surfaces every
distinct failure mode we have observed.

\bibliographystyle{IEEEtran}
\bibliography{../../bibliography/refs}

\end{document}
